\emph{(i)} Write $\dot\nabla$ for the derivative acting on the dotted factor only. For a bivector $F$ and a vector $A$, $\nabla\cdot(F\cdot A) = (\nabla\cdot F)\cdot A + \langle F\dot A\dot\nabla\rangle_{0}$, and $\langle F\dot A\dot\nabla\rangle_{0} = F\cdot(\dot A\wedge\dot\nabla) = -F\cdot(\nabla\wedge A)$. With $F_{a} = \nabla\wedge A_{a}$ this gives
\[
\nabla\cdot\omega = (\nabla\cdot F_{1})\cdot A_{2} - (\nabla\cdot F_{2})\cdot A_{1} - F_{1}\cdot F_{2} + F_{2}\cdot F_{1},
\]
and the scalar product of bivectors is symmetric, so the last two terms cancel. Since $\nabla\wedge F = \nabla\wedge\nabla\wedge A = 0$, the free equation $\nabla F = 0$ is $\nabla\cdot F = 0$, on which $\nabla\cdot\omega = 0$. The check verifies the identity for two generic polynomial potentials of degree two with $\nabla\cdot F_{1} \ne 0$ and $F_{1}\cdot F_{2} \ne 0$, so neither side is trivially zero.

\emph{(ii)} $n\cdot(n\cdot F) = (n\wedge n)\cdot F = 0$, so $n\cdot F$ has no component along $\bar n$ (as $n\cdot\bar n = 2 \ne 0$) and is $\kappa n + E_{\perp}$. Under $A \mapsto A + \nabla\chi$, $F$ is unchanged ($\nabla\wedge\nabla\chi = 0$), hence so is $E_{\perp}$. For the energy, decompose $F = \alpha\,n\wedge\bar n + n\wedge p + \bar n\wedge r + \beta J$ with $p$, $r$ transverse. Then $n\cdot(n\wedge\bar n) = -2n$, $n\cdot(n\wedge p) = 0$, $n\cdot(\bar n\wedge r) = 2r$, $n\cdot J = 0$, so $n\cdot F = -2\alpha n + 2r$ and $E_{\perp} = 2r$. Using $Q(n) = 0$, $n\bar nn = 4n$ and that transverse vectors anticommute with $n$ and $\bar n$ while $J$ commutes with them, $n(n\wedge\bar n)n = n(n\bar n - 2)n = 0$, $n(n\wedge p)n = 0$, $nJn = 0$ and $n(\bar n\wedge r)n = n\bar nrn = -n\bar nnr = -4nr$; so $nFn = -4nr$ and $\langle nFnF\rangle_{0} = -4\langle nrF\rangle_{0}$. Of the four pieces only $\bar n\wedge r$ contributes, $\langle nr\bar nr\rangle_{0} = -\langle n\bar nrr\rangle_{0} = -2Q(r)$ (the others give $\langle nrn\bar n\rangle_{0} = 0$, $\langle nrnp\rangle_{0} = 0$, $n\cdot(rJ) = 0$), hence $n\cdot T(n) = -\tfrac12\langle nFnF\rangle_{0} = -4Q(r) = -E_{\perp}^{2}$, positive when $E_{\perp} \ne 0$ since $E_{\perp}$ is spacelike (checked on three generic bivectors). In the gauge $n\cdot A = 0$, $n\cdot F = n\cdot(\nabla\wedge A) = \partial_{n}A - \nabla(n\cdot A) = \partial_{n}A$. For a vector $X$, $n\cdot X = \langle nX\rangle_{0}$, and $n\cdot(F\cdot A) = \langle nFA\rangle_{0} = (n\cdot F)\cdot A$ because $\langle (n\wedge F)A\rangle_{0} = 0$. So $n\cdot\omega = \partial_{n}A_{1}\cdot A_{2} - \partial_{n}A_{2}\cdot A_{1}$; with $A_{c} = \mu_{c}n + A_{c\perp}$, $Q(n) = 0$ and $n\cdot A_{c\perp} = 0$ the coefficients $\mu_{c}$ drop out, leaving $\partial_{n}A_{1\perp}\cdot A_{2\perp} - \partial_{n}A_{2\perp}\cdot A_{1\perp}$ (checked identically as polynomials, for potentials with nonzero $\mu_{c}$). The flux density contains no transverse derivative, so the form $\int_{N}n\cdot\omega$ is a sum over generators of forms in the data of each generator alone; brackets of data on distinct generators vanish.

\emph{Modes.} On one generator closed to period $2\pi$, with a transverse cell of area $s$ and $A_{i}$ the components of $A_{\perp} = A_{1}\gamma_{1} + A_{2}\gamma_{2}$, $A_{\perp}\cdot A_{\perp}' = -\sum_{i}A_{i}A_{i}'$. By orthogonality of $\cos k\lambda$, $\sin k\lambda$ on $[0, 2\pi]$ (each squared integrating to $\pi$), $\int_{0}^{2\pi}\partial_{\lambda}A_{i}\,A_{i}'\,d\lambda = \pi\sum_{k}k(b_{ik}a_{ik}' - a_{ik}b_{ik}')$, so the flux of $\omega$ is $\varpi = -s\sum_{i}\int(\partial A_{i}A_{i}' - \partial A_{i}'A_{i}) = 2\pi s\sum_{i,k}k(a_{ik}b_{ik}' - b_{ik}a_{ik}')$, and $P = s\int -E_{\perp}^{2} = s\sum_{i}\int(\partial_{\lambda}A_{i})^{2} = \pi s\sum k^{2}(a_{ik}^{2} + b_{ik}^{2})$. The null translation $\lambda \mapsto \lambda + t$ moves $(a_{ik}, b_{ik})$ with velocity $(kb_{ik}, -ka_{ik})$, and $\varpi(X, Y) = dP(Y)$ for every $Y$ exactly for $X$ this velocity: $2\pi sk(X_{a}Y_{b} - X_{b}Y_{a}) = 2\pi sk^{2}(aY_{a} + bY_{b})$ forces $X_{a} = kb$, $X_{b} = -ka$. So $P$ generates the future null translations, and the corresponding bracket is $\{a_{ik}, b_{jl}\} = \delta_{ij}\delta_{kl}/(2\pi sk)$ (then $\{a_{ik}, P\} = kb_{ik}$). With $[\,,\,] = J\{\,,\,\}$ ($J$ central in the algebra of the modes, $J^{*} = -J$, $a$, $b$ self-adjoint), $[a_{ik} + Jb_{ik}, a_{ik} - Jb_{ik}] = -2J[a_{ik}, b_{ik}] = 1/(\pi sk)$, so $[\alpha_{ik}, \alpha_{ik}^{*}] = k$; modes with different $(i, k)$ commute, $[\alpha_{ik}, \alpha_{jl}] = (\pi s)kl([a, a] - [b, b] + J[a, b] + J[b, a]) = 0$; and $[P, a + Jb] = \pi sk^{2}(2J a\,c - 2b\,c)$ with $c = J/(2\pi sk)$, which is $-k(a + Jb)$. Replacing $J$ by $-J$ in both the quantisation rule and the modes ($\alpha_{ik}' = (\pi s)^{1/2}k(a_{ik} - Jb_{ik})$) gives $[\alpha_{ik}', \alpha_{ik}'^{*}] = -2(-J)[a_{ik}, b_{ik}]' \cdot \pi sk^{2}$ with $[a, b]' = -J/(2\pi sk)$, again $+k$: the sign is fixed by the bracket, that is by positive null energy and the future direction, and not by the orientation of the unit.

\emph{(iii)} Both sides are linear in the transverse vector $u$. For $u = \gamma_{1}$: $I_{4}n\gamma_{1} = -\gamma_{0}\gamma_{2} + \gamma_{2}\gamma_{3}$ and $n\gamma_{1}J = (\gamma_{0}\gamma_{1} + \gamma_{3}\gamma_{1})\gamma_{1}\gamma_{2} = -\gamma_{0}\gamma_{2} + \gamma_{2}\gamma_{3}$; for $u = \gamma_{2}$: $I_{4}n\gamma_{2} = \gamma_{0}\gamma_{1} - \gamma_{1}\gamma_{3} = n\gamma_{2}J$. So $I_{4}(n\wedge u) = (n\wedge u)J = n\wedge(uJ)$, with $uJ$ the quarter turn of $u$ in the transverse plane, and $J(n\wedge u) = n(Ju) = -n(uJ)$ because $J$ commutes with $n$ and anticommutes with $u$ (checked also for $u = 3\gamma_{1} - 2\gamma_{2}$). For $\varepsilon = \gamma_{1}E_{\perp} = -E_{1} + E_{2}J \in \mathrm{span}\{1, J\}$ and $R = e^{-J\phi/2}$, $RE_{\perp}\tilde R = E_{\perp}\tilde R^{2}$ (transverse vectors anticommute with $J$), so $\gamma_{1}RE_{\perp}\tilde R = \varepsilon e^{J\phi}$ (checked with the unnormalised rotor $2 + J$). In components the rotation is $(E_{1}, E_{2}) \mapsto (\cos\phi\,E_{1} + \sin\phi\,E_{2}, -\sin\phi\,E_{1} + \cos\phi\,E_{2})$, and on the modes, which are linear in the components, $\alpha_{1k} \mp J\alpha_{2k} \mapsto (\cos\phi \pm J\sin\phi)(\alpha_{1k} \mp J\alpha_{2k})$; by (ii) both evolve by $e^{-Jkt}$ under the null translation.

\emph{(iv)} Write $a, b, c, d$ for the labels $x, y, z, u$ (a generator and a channel each). The CAR identity $[v_{a}w_{b}, v_{c}w_{d}] = \delta_{bc}v_{a}w_{d} - \delta_{ad}v_{c}w_{b}$ (distinct planes anticommute, $\{v_{j}, w_{j}\} = 1$, Fact \claimref{field-theory.imported-tower-limit}) summed over the window gives the displayed commutator, the first sum over $k$ with $k$, $k + m$, $k + m + l$ in the window and the second over $k$ with $k$, $k + l$, $k + l + m$ there. Each sum differs from the full current $j^{ad}_{m+l}$, respectively $j^{cb}_{m+l}$, by bilinears $c_{x,k}^{*}c_{y,k'}$ with $k$ or $k'$ within $\max(|m|, |l|)$ of an end of the window. On a vector edge-clean at that depth such a bilinear vanishes unless it is a number operator at the bottom (where $c^{*}$ hits an occupied level otherwise, or $c$ at the top hits an empty one); the number operators occur only for $a = d$, $b = c$, $l = -m$, and contribute $\sum_{k \le m}n_{a,k} = m$ (bottom levels filled), so the commutator is $\delta_{bc}j^{ad}_{m+l} - \delta_{ad}j^{cb}_{m+l} + m\,\delta_{ad}\delta_{bc}\delta_{m+l,0}$ there. For $a = b = c = d$, $l = -m$ this is Theorem \claimref{hawking.current-commutator-exactly-edge-split}. The check verifies the exact identity for all $a, b, c, d$ and $|m|, |l| \le 1$ at $L = 3$, and $[j^{ab}_{1}, j^{ba}_{-1}] = 1 + j^{aa}_{0} - j^{bb}_{0}$ on the four edge-clean basis vectors. $[n_{a,k}, c_{a,k}^{*}c_{b,k+m}] = c_{a,k}^{*}c_{b,k+m}$ and $[n_{b,k+m}, c_{a,k}^{*}c_{b,k+m}] = -c_{a,k}^{*}c_{b,k+m}$ give $[H, j^{ab}_{m}] = (k - (k + m))j^{ab}_{m} = -mj^{ab}_{m}$. On a null plane the generator CAR is $\{\psi_{N}(\lambda, y), \psi_{N}^{*}(\lambda', y')\} = \delta(\lambda - \lambda')\delta^{2}(y - y')\Lambda(\bar n)$ (\claimref{synthesis.bulk-hadamard-states-restrict-to-the-generator-condition-on-flat-null-planes}(i)), so modes of distinct generators are distinct planes and their pair currents, being even, commute.

\emph{(iv), the rotation.} For one mode, $(1 + (z - 1)n)\,c\,(1 + (\bar z - 1)n) = \bar z c$ and $c^{*} \mapsto zc^{*}$, since $nc = 0$, $cn = c$ and $J$ commutes with every transport generator. Hence $c_{a,k}^{*}c_{b,k+m} \mapsto z_{a}\bar z_{b}\,c_{a,k}^{*}c_{b,k+m}$: the densities ($a = b$) are invariant, $j^{+-}_{m}$ is multiplied by $z_{+}^{2}$ and $j^{-+}_{m}$ by $z_{-}^{2}$. $U^{*}U = 1$ because each factor has adjoint $1 + (\bar z - 1)n$ ($J^{*} = -J$, $n^{*} = n$) and $z\bar z = 1$. The phases $z_{\pm} = e^{\pm J\phi/2}$ are those of the two channels under the rotation about the generator in \claimref{synthesis.light-ray-good-part-of-a-spinor-is-two-chirality-arrays}(iv), read with the tower's unit $J$; the check uses $z_{\pm} = (3 \pm 4J)/5$. By the identity at $x = y$, $z = u$ the densities of distinct channels commute and $[j^{xx}_{m}, j^{xx}_{-m}] = m$ on edge-clean vectors: exact canonical pairs, invariant under $U(\phi)$.

\emph{(v)} The span of the pair currents of one generator over $\mathrm{span}\{1, J\}$ decomposes under $U(\phi)$ into the phases $1$ (densities), $z_{+}^{2}$ ($j^{+-}_{m}$) and $z_{-}^{2}$ ($j^{-+}_{m}$), so its helicity $+1$ part is spanned by the $j^{+-}_{m}$. By (iv), $[j^{+-}_{m}, j^{+-}_{l}] = 0$ ($\delta_{bc} = \delta_{ad} = 0$) and on edge-clean vectors $[j^{+-}_{m}, j^{-+}_{l}] = j^{++}_{m+l} - j^{--}_{m+l} + m\,\delta_{m+l,0} = 2j^{3}_{m+l} + m\,\delta_{m+l,0}$. Let $\Omega$ be the bottom-filled window (each channel filled up to a level $h$ with the depth conditions); by \claimref{synthesis.fluctuation-forms-of-the-tower-and-density-wave-bosons-of-the-null-vacuum}(vi) this is the restriction to band modes of the null vacuum of \claimref{synthesis.null-vacuum-on-a-horizon-generator}(ii), the filled sea of negative null momentum on every generator. For $q < 0$, $j^{aa}_{q}\Omega = \sum_{k}c_{a,k}^{*}c_{a,k+q}\Omega$ is a sum of $|q|$ orthogonal unit vectors (a particle at $k > h$ and a hole at $k + q \le h$), so $\|j^{aa}_{q}\Omega\|^{2} = |q|$, and the two channels give orthogonal vectors: $\|2j^{3}_{q}\Omega\|^{2} = 2|q|$. (Checked: $2$ for $(m, l) = (0, -1)$ at $L = 3$ and $(1, -2)$ at $L = 4$, and $4$ for $(0, -2)$ and $(-1, -1)$ at $L = 4$.)

\emph{Many generators.} With $g, h$ labelling generators, $\varepsilon_{m}^{*}$ has coefficients $\bar M$ and the roles of the channels exchanged, and the CAR identity gives, on vectors edge-clean in every generator,
\[
[\varepsilon_{m}, \varepsilon_{l}^{*}] = \sum_{g,g'}(MM^{*})_{gg'}j^{g+,g'+}_{m-l} - \sum_{h,h'}(M^{*}M)_{h'h}j^{h-,h'-}_{m-l} + m\,\delta_{ml}\operatorname{tr}MM^{*},
\]
with $j^{x,y}_{q} = \sum_{k}c_{x,k}^{*}c_{y,k+q}$ across generators. For $m = l$ the operator part has expectation $\sum_{g}(MM^{*})_{gg}h - \sum_{h}(M^{*}M)_{hh}h = 0$ on $\Omega$, so the central value is $m\operatorname{tr}MM^{*}$. For $m < l$, $q = m - l < 0$, the vectors $j^{x,y}_{q}\Omega$ for distinct pairs $(x, y)$ are orthogonal (distinct occupation patterns) and each has norm squared $|q|$, as in the one-generator case (for $x \ne y$ the particle and the hole sit in different channels, with the same count of positions). So $\|D\Omega\|^{2} = |q|(\|MM^{*}\|_{F}^{2} + \|M^{*}M\|_{F}^{2}) = 2|q|\,\|MM^{*}\|_{F}^{2}$, since $\|M^{*}M\|_{F}^{2} = \operatorname{tr}(M^{*}MM^{*}M) = \operatorname{tr}(MM^{*}MM^{*})$. For the self-adjoint $H = MM^{*}$ on $N$ generators, Cauchy and Schwarz give $(\operatorname{tr}H)^{2} \le N\operatorname{tr}H^{2} = N\|H\|_{F}^{2}$, with equality exactly when $H$ is scalar. (Checked on two generators for $M = 1$, $\begin{pmatrix}1&1\\0&1\end{pmatrix}$, $\begin{pmatrix}2&0\\1&1\end{pmatrix}$: norms $4$, $14$, $56$, central values $2$, $3$, $6$, equality only for $M = 1$; on three generators with $M = 1$: norm $6$, central value $3$.) For $M = \mathrm{diag}(s^{1/2}f(y_{g}))$, $\operatorname{tr}MM^{*} = s\sum_{g}|f(y_{g})|^{2}$ and $\|MM^{*}\|_{F}^{2} = s^{2}\sum_{g}|f(y_{g})|^{4}$, a Riemann sum for $s\int|f|^{4}d^{2}y$, which tends to $0$ as $s \to 0$ for bounded $f$ of compact support (argued). The same bound applied to $\varepsilon_{\pm}$ built from $j^{+-}$ and $j^{-+}$ gives the relations stated, the commutator $[\alpha_{+,k}, \alpha_{-,l}^{*}]$ corresponding to $[j^{+-}_{k}, j^{+-}_{-l}] = 0$ exactly and $[\alpha_{+,k}, \alpha_{-,l}]$ to $2j^{3}_{k+l}$, which annihilates $\Omega$ for $k, l \ge 1$.
